% Recommended papers: Tools and datasets for mining libre software repositories
%			The promises and perils of mining git
%
%
%
%
%

% $Id: slides.tex 4228 2006-06-21 21:55:12Z jjamor $
%
%
% Compilar a .pdf con LaTeX (pdflatex)
% Es necesario instalar Beamer (paquete latex-beamer en Debian)
%

%
% Gráficos:
% Los gráficos pueden suministrarse en PNG, JPG, TIF, PDF, MPS
% Los EPS deben convertirse a PDF (usar epstopdf)
%

\documentclass{beamer}
\usetheme{Warsaw}
\usebackgroundtemplate{\includegraphics[width=\paperwidth]{format/libresoft-bg.png}}
\usepackage[spanish]{babel}
\usepackage[utf8]{inputenc}
\usepackage{graphics}
\usepackage{amssymb} % Simbolos matematicos
\usepackage{url}

%\definecolor{libresoftgreen}{RGB}{162,190,43}
%\definecolor{libresoftblue}{RGB}{0,98,143}

%\setbeamercolor{titlelike}{bg=libresoftgreen}

%% Metadatos del PDF.
\hypersetup{
  pdftitle={FLOSS Data Sources},
  pdfauthor={Daniel Izquierdo},
  pdfcreator={GSyC/Libresoft},
  pdfproducer=PDFLaTeX,
  pdfsubject={nn},
}
%%


\AtBeginSection[]
{
  \begin{frame}<presentation>
    \frametitle{Index}
    \tableofcontents[current]
  \end{frame}
}


\begin{document}

\title{Introduction to FLOSS Data Sources}
\subtitle{Master on Free Software}
\institute{\\dizquierdo@libresoft.es\\
GSyC/Libresoft}
\author{Daniel Izquierdo}
\date{\today}

\frame{
\maketitle
\begin{center}
\includegraphics[width=6cm]{format/gsyc-urjc}
\end{center}
}


% Si el titulo o el autor se quieren acortar para los pies de página
% se pueden redefinir aquí:
%\title{Titulo corto}
%\author{Autores abreviado}


%% LICENCIA DE REDISTRIBUCION DE LAS TRANSPAS
\frame{
~
\vspace{4cm}

\begin{flushright}
{\tiny
(cc) 2011 Daniel Izquierdo. \\
Some rights reserved. This document is distributed under the Creative \\
            Commons Attribution-ShareAlike 3.0 licence, available in \\
            http://creativecommons.org/licenses/by-sa/3.0/

%  Este documento (o uno muy similar) está disponible en \\
%  \url{http://gsyc.escet.urjc.es/~jjamor/}
}
\end{flushright}
}
%%

%%%%%%
%Transpas separadas por \begin{frame}
%%%%%%%%%%%%%%%%%%%%%%%%\end{frame}

\section{Introduction}

\begin{frame}
\frametitle{}
\begin{itemize}
\item Global software development implies a distributed environment of work
\item In the specific case of FLOSS projects, specific infrastructure is 
used as communication channels.
\end{itemize}
\end{frame}

%%%%%%%%%%%%%%%%%%%%%%%%%%%%%%%%%%%%%%%%%%%%%%%%%%%%%%%%%%%%%%

\begin{frame}
 \frametitle{}
 \begin{itemize}
 \item That infrastructure is generally speaking, publicly available
 \item Thanks to this open nature of the repositories of information, it 
 is possible to study them
 \end{itemize}
\end{frame}

%%%%%%%%%%%%%%%%%%%%%%%%%%%%%%%%%%%%%%%%%%%%%%%%%%%%%%%%%%%%%%


\section{Data Sources}
\begin{frame}
 \frametitle{Main data sources}
 \begin{itemize}
  \item Source code management system
  \item Mailing lists
  \item Bug tracking system
  \item Source code
 \end{itemize}
\end{frame}

%%%%%%%%%%%%%%%%%%%%%%%%%%%%%%%%%%%%%%%%%%%%%%%%%%%%%%%%%%%%%%

\begin{frame}
 \frametitle{Main data sources: SCM}
 \begin{itemize}
 \item This is used to manage file versions during the development process
 \item It is possible to obtain the state of the source code at any time during its history
 \end{itemize}
\end{frame}



\begin{frame}
 \frametitle{Main data sources: SCM}
 \begin{itemize}
 \item Its main pieces of information are the following:
  \begin{itemize}
  \item Who:  Author/Committer
  \item When: Date of commit
  \item What: Files and lines \emph{touched}
  \end{itemize}
 \end{itemize}


\end{frame}



\begin{frame}
 \frametitle{Main data sources: SCM}
 \begin{itemize}
 \item There are two main types: centralized and distributed
 \item Centralized: CVS or SVN
 \item Distributed: Git or Mercurial
 \end{itemize}
\end{frame}



\begin{frame}
 \frametitle{Main data sources: SCM}
 \begin{itemize}
 \item \emph{git log | hg log | svn log ...}
 \item \emph{git diff | hg diff | ...}
 \end{itemize}
\end{frame}



\begin{frame}
 \frametitle{Main data sources: Mailing lists}
 \begin{itemize}
 \item Used to organized libre software projects
 \item Asynchronous way of communication
 \item Basically, those forward received e-mail messages to subscribed e-mail addresses
 \item Usually stored in specific formats such as RFC 4155

 \end{itemize}

\end{frame}


\begin{frame}
 \frametitle{Main data sources: Mailing lists}
 \begin{itemize}
 \item Some interesting headers:
  \begin{itemize}
   \item From: this field indicates the sender of the e-mail
   \item To: this field indicates the receiver of the e-mail
   \item Message-ID: unique identifier of the email
   \item In-Reply-To: indicates who is the \emph{father} of this e-mail (will refer to the Message-ID field)
   \item List-id: distribution list
  \end{itemize}

 \end{itemize}


\end{frame}

\begin{frame}
 \frametitle{Main data sources: Mailing lists}
 \begin{itemize}
 \item This type of tool is also used as a peer review process of the changes done in the SCM
   \begin{itemize}
     \item Lists of Commits-watchers
     \item Patches attached 
   \end{itemize}

 \end{itemize}


\end{frame}

\begin{frame}
 \frametitle{Main data sources: Bug tracking system}
 \begin{itemize}
 \item This tool is used to manage the evolution of the several errors detected in the source code


 \end{itemize}


\end{frame}

\begin{frame}
 \frametitle{Main data sources: BTS}
 \begin{itemize}
 \item Some examples:
   \begin{itemize}
    \item Bugzilla
    \item Sourceforge tracker
    \item GNATS: used by the FSF
    \item Jira
    \item Mantis
   \end{itemize} 


 \end{itemize}
\end{frame}

\begin{frame}
 \frametitle{Main data sources: BTS}
 \begin{itemize}
 \item Life cycle: \url{http://www.bugzilla.org/docs/2.18/html/lifecycle.html}

 \end{itemize}

\end{frame}

\begin{frame}
 \frametitle{Main data sources: BTS}
 \begin{itemize}
 \item Some interesting fields:
   \begin{itemize}
    \item Bug id: unique identifier
    \item Description
    \item Opened: date when the bug report was sent
    \item Status: current status of the report - new, assigned, reopened, needinfo, verified, etc...
    \item Resolution: invalid, wontfix, notabug, fixed, and others...
 \end{itemize}

 \end{itemize}
\end{frame}

\begin{frame}
 \frametitle{Main data sources: Source code}
 \begin{itemize}
 \item Generally obtained by means of the releases of source code
 \item bug also from the trunk or master branch of the SCMs

 \end{itemize}

\end{frame}

\begin{frame}
 \frametitle{Main data sources: Source code}
 \begin{itemize}
 \item There exist a hierarchical structure
 \item It is possible to discriminate by type of file (translation, source code, binary, makefile, etc).
 \item Definitively, several types of files, and not all of them are useful for evaluating software!

 \end{itemize}

\end{frame}



\begin{frame}
 \frametitle{Main data sources: Source code}
 \begin{itemize}
 \item Also interesting to evaluate the copyright and authors of 
  the several files
 \item This makes possible similar studies about ownership in the source code as done in the SCM
 \end{itemize}

\end{frame}


\begin{frame}
 \frametitle{Main data sources: Source code}
 \begin{itemize}
 \item Also interesting to study licenses (in the end, the license is found
 in the source code).
 \item Between the source code and what a developer claims, the source code will 
 provide the actual information.
 \end{itemize}

\end{frame}

\section{References}

\begin{frame}
 \frametitle{References}
 \begin{itemize}
 \item Producing OSS by \textit{Karl Fogel}
 \item Tools and datasets for mining libre software repositories, by
 \textit{Gregorio Robles, Jes\'us M. Gonz\'alez-Barahona, Daniel Izquierdo-Cort\'azar and Israel Herraiz}
 \item The promises and perils of mining git by
 \textit{Christian Bird, Peter Rigby, Earl Barr, David Hamilton, Daniel German, and Premkumar Devanbu}
 \end{itemize}

\end{frame}

\end{document}